% ============================================================
%  Chapitre 1 — Introduction Générale
% ============================================================
\chapter{Introduction Générale}
\thispagestyle{fancy}

\begin{tcolorbox}[technote, title={Épigraphe}]
\itshape
« L'intelligence artificielle est la nouvelle électricité.\\
Tout comme l'électricité a transformé de nombreuses industries il y a cent ans,\\
l'IA va maintenant transformer presque toutes les industries du monde. »
\begin{flushright}
\upshape— \textbf{Andrew Ng}, Co-fondateur de Coursera, ex-directeur de l'IA chez Baidu
\end{flushright}
\end{tcolorbox}

\section{Contexte général}

L'agriculture demeure le pilier central de l'économie tunisienne, représentant
\textbf{10\% du PIB} et \textbf{14\% de l'emploi national} selon les données de
l'Institut National de la Statistique \cite{ins_tunisie_2023}. Ce secteur nourrit plus de
11 millions de Tunisiens et génère des exportations significatives — huile d'olive,
dattes, agrumes, poissons — vers les marchés européens et arabes.

Pourtant, malgré son importance stratégique, l'agriculture tunisienne fait face à une
convergence de défis structurels majeurs :

\begin{itemize}
  \item \textbf{Stress hydrique croissant :} Avec une pluviométrie moyenne inférieure à
        400 mm/an dans les régions intérieures et une surexploitation des nappes
        phréatiques, la gestion de l'eau est devenue critique.

  \item \textbf{Pression des ravageurs et maladies :} Les pertes agricoles attribuées aux
        maladies fongiques, bactériennes et aux insectes nuisibles représentent 20 à 40\%
        des récoltes dans certaines régions, selon la FAO \cite{fao_tunisie_2022}.

  \item \textbf{Déclin des pollinisateurs :} Le secteur apicole tunisien, avec ses
        300 000 ruches environ, connaît un taux de mortalité hivernal alarmant lié au
        Varroa destructor, aux pesticides et aux changements climatiques.

  \item \textbf{Vieillissement des exploitants :} L'âge moyen des agriculteurs tunisiens
        dépasse 50 ans, avec un déficit de transmission des savoirs et une faible adoption
        des technologies numériques.

  \item \textbf{Fragmentation des exploitations :} Plus de 70\% des fermes couvrent moins
        de 5 hectares, limitant les économies d'échelle et l'accès aux équipements modernes.
\end{itemize}

\section{L'Agriculture 4.0 — opportunités et enjeux}

Face à ces défis, l'émergence de l'\textbf{Agriculture 4.0} — convergence de l'IoT, de
l'intelligence artificielle, des drones et de l'analyse big data — ouvre des perspectives
transformatrices. Les projections mondiales indiquent que le marché de l'agritech basée
sur l'IA devrait atteindre \textbf{4 milliards de dollars en 2026}, avec un taux de
croissance annuel de 25\% \cite{kamilaris2018deep}.

Les technologies clés de cette révolution agricole incluent :

\begin{itemize}
  \item \textbf{Capteurs IoT et télémétrie :} Surveillance en temps réel de l'humidité du
        sol, de la température, du débit d'eau et de l'état des cultures.
  \item \textbf{Vision par ordinateur :} Détection automatique des maladies foliaires par
        analyse d'images, avec des précisions supérieures à 90\% pour les modèles YOLO
        récents \cite{ultralytics2023yolov8}.
  \item \textbf{Grands modèles de langage :} Assistants agricoles conversationnels capables
        de répondre en langue locale aux questions des exploitants.
  \item \textbf{MLOps :} Industrialisation des modèles d'IA avec suivi des expériences,
        versionnement des données et déploiement continu \cite{mlflow2023}.
\end{itemize}

\section{Présentation du projet \farmAI}

\farmAI{} v3.0 est une plateforme agro-technologique complète, conçue pour démocratiser
l'accès aux technologies d'IA dans le milieu agricole tunisien. Le projet se distingue
par quatre piliers fondamentaux :

\begin{tcolorbox}[infobox, title={Les 4 piliers de Smart Farm AI}]
\begin{enumerate}
  \item \textbf{IoT Temps Réel :} Deux nœuds ESP32 (NODE\_A irrigation, NODE\_B rucher)
        transmettant des données télémétriques toutes les 10 secondes via HTTP POST.
  \item \textbf{Vision IA Souveraine :} Huit modèles YOLO v11 couvrant 13 espèces et
        pathologies agricoles, avec inférence locale sans dépendance cloud.
  \item \textbf{Assistant Darija :} LLM Labess-7B spécialisé en dialecte tunisien,
        enrichi par RAG avec des sources officielles (UTAP, AVFA), fonctionnant
        entièrement en local via Ollama.
  \item \textbf{MLOps Industriel :} Pipeline DVC + MLflow avec versionnement des données,
        entraînement reproductible et registre de modèles.
\end{enumerate}
\end{tcolorbox}

La plateforme adopte une architecture \textbf{multi-tenant} permettant à plusieurs
exploitations agricoles de partager la même infrastructure tout en maintenant
l'isolation complète de leurs données. Elle est accessible depuis une interface web
\textbf{Progressive Web App} (PWA) utilisable hors-ligne par les ouvriers agricoles,
et depuis un tableau de bord avancé pour les propriétaires d'exploitations.

\section{Problématique}

La problématique centrale de ce projet peut être formulée ainsi :

\begin{tcolorbox}[warningbox, title={Problématique}]
\textit{Comment concevoir et déployer une plateforme d'intelligence artificielle
accessible, souveraine et opérationnellement robuste pour les agriculteurs tunisiens,
intégrant surveillance IoT en temps réel, détection automatique des pathologies agricoles
par vision par ordinateur, et assistance conversationnelle en dialecte local — tout en
garantissant la reproductibilité des modèles ML via des pratiques MLOps industrielles
et la souveraineté technologique par l'exécution locale des modèles ?}
\end{tcolorbox}

\section{Objectifs du projet}

Les objectifs spécifiques poursuivis sont :

\begin{enumerate}
  \item Concevoir une architecture système évolutive (FastAPI + React + PostgreSQL) supportant
        la gestion multi-fermes et multi-utilisateurs.
  \item Développer et déployer 8 modèles YOLO v11 pour la détection de pathologies
        agricoles et la surveillance des ruches.
  \item Implémenter un système IoT complet avec deux nœuds ESP32, simulés via Wokwi et
        opérationnels sur matériel réel.
  \item Créer un assistant conversationnel souverain en Darija tunisien, utilisant le LLM
        Labess-7B et un système RAG avec ChromaDB.
  \item Mettre en place une chaîne MLOps complète (DVC + MLflow) pour la reproductibilité
        et le suivi des modèles en production.
  \item Développer un module de gestion apicole complet (HiveIntel) avec 7 sous-modules
        fonctionnels.
  \item Déployer la solution sur Oracle Cloud Free Tier en mode haute disponibilité.
\end{enumerate}

\section{Plan du rapport}

Ce rapport s'articule en sept chapitres, organisés comme suit :

\begin{description}
  \item[\textbf{Chapitre 1 — Introduction générale :}] Contexte, problématique et
        objectifs du projet (présent chapitre).

  \item[\textbf{Chapitre 2 — Contexte et étude préliminaire :}] Analyse du secteur
        agricole tunisien, benchmark des solutions existantes, cahier des charges
        fonctionnel et non-fonctionnel.

  \item[\textbf{Chapitre 3 — Méthodologies :}] Présentation des méthodologies adoptées :
        Scrum (6 sprints sur 14 semaines) et CRISP-DM pour la conduite du projet Data
        Science.

  \item[\textbf{Chapitre 4 — Architecture et conception :}] Diagrammes UML (cas
        d'utilisation, classes, séquences), architecture système en couches et
        architecture IoT.

  \item[\textbf{Chapitre 5 — MLOps et pipeline IA :}] Détail de la chaîne MLOps,
        présentation des 8 modèles YOLO, module IsolationForest et assistant souverain
        RAG-Darija.

  \item[\textbf{Chapitre 6 — Réalisation et tests :}] Environnement de développement,
        présentation des interfaces utilisateur, tableau des endpoints API et résultats
        de tests.

  \item[\textbf{Chapitre 7 — Conclusion et perspectives :}] Synthèse des réalisations,
        difficultés rencontrées et perspectives d'évolution.
\end{description}
